\section[Numérisation à Cujas]{Les balbutiements de la numérisation à Cujas (2001 – 2007)}

Le choix de la bibliothèque Cujas est d'orienter d'emblée la valorisation de son contenu vers \enquote{un lectorat universitaire de haut niveau}. Noëlle Balley et Sébastien Dalmon attribue cette volonté à la spécialisation extrême du fonds: à leurs yeux, ce contenu \enquote{souvent aride} s'avère peu proprice à une \enquote{valorisation orientée vers le grand public}. L'établissement privilégie donc des corpus spécifiquement sélectionnés pour ce public académique, en mettant à sa disposition des \enquote{outils facilitant leur travail sur les textes}. Toutefois, ils soulignent que malgré les moyens mis en oeuvre pour s'adresser à ce lectorat élitiste, il convient de ne pas \enquote{négliger le grand public}.

L'un des premiers chantiers de numérisation de la bibliothèque Cujas remonte aux années 2000 et 2001. Durant cette phase pionnière, l'établissement fait numériser 17 titres représentant environ 22 volumes, une tâche externalisée auprès de la société Safig. La bibliothèque commande alors une numérisation \enquote{en mode image, avec une navigation grâce à la pagination}. Séduite par l'offre du prestataire qui propose de prendre en charge l'hébergement et la diffusion en ligne des documents sur ses propres serveurs, la bibliothèque Cujas choisit de lui déléguer l'ensebmle de cette chaîne de stockage et de mise à disposition. 

Cette stratégie se révèle désastreuse : la société victime d'une attaque de ses serveurs, entraînant la perte irrémédiable d'une majeure partie des documents numérisés. Face à ce sinistre, la bibliothèque s'efforce de récupérer ce qu'elle peut. Le prestataire Safig restitue ce qui lui reste \enquote{sous forme de CD-ROM}. A la réception des CD-ROM, le personnel de Cujas tente de déterminer s'il est possible \enquote{de réutiliser du contenu}. Cette tentative de sauvetage s'avère toutefois vaine : les fichiers réceptionnés se révèlent irrécupérables, voire corrompus pour certains, contraignant l'établissement à prendre la décision de renumériser intégralement le corpus.

Le 5 et 6 novembre 2008, la bibliothèque Cujas a organisé un \enquote{colloque international en collaboration avec le Sénat},portant sur le doyen Jean Carbonnier. Noëlle Balley et Sébastien Dalmon qualifient cette initiative de \enquote{première entreprise de numérisation effective [qui] a concerné, pour des raisons conjoncturelles, les polycopiés du doyen Jean Carbonnier}. La numérisation de ces polycopiés a été réalisé à \enquote{l'occasion du centenaire de la naissance de ce grand juriste civiliste}. A l'origine, ce projet résidait dans une \enquote{participation modeste} de la bibliothèque Cujas à \enquote{un projet éditorial}, initié par l'anthropologue du droit Raymond Verdier. Son objectif principal était alors \enquote{de rassembler les écrits de Jean Carbonnier introuvables, inédits ou dispersés}. La valorisation de ce projet scientifique a revêtu diverses formes, comme la tenue de colloques au Sénat, des actions de commmunication, une exposition virtuelle.

Parallèlement à la tenue des colloques scientifiques, la bibliothèque Cujas prend l'initiative d'organiser une exposition virtuelle d'envergure. L'établissement fait ainsi le choix \enquote{de mettre en ligne une exposition virtuelle consacrée à Jean Carbonnier, constatant l'absence, sur Internet, de ressources suffisamment fiables et complètes sur le Doyen}. L'objectif est à la fois de faire connaitre \enquote{une grande figure du droit contemporain, méconnue du grand public, mais aussi à répondre aux attentes des spécialistes}. 

Pour mener à bien cette valorisation cutlurelle, la bibliothèque Cujas s'appuie opportunément sur son expérience aquise en matière d'éditorialisation numérique. L'établissement n'est en effet pas à son coup d'essai : il avait déjà réalisé une expostion sur\enquote{le Bicentenaire du Code civil (1804-2004) en collaboration avec la Cour de cassation et l'Ordre des avocats}. Cette antériorité partenariale a permis aux équipes de Cujas de mmobiliser un savoir-faire technique pour concevoir une exposition virtuelle, entièrement dédié à la mise en valeur scientifique et mémorielle de jean Carbonnier.

L'exposition virtuelle consacrée à Jean Carbonnier offre un \enquote{accès à treize polycopiés de cours du doyen Carbonnier professés entre 1958 et 1974, portant sur la sociologie juridique ou le droit civil}. Ce corpus de documents, qui constitue la première collection intégrée à la bibliothèque numériuque, est aujourd'hui hébergé sur la plateforme actuelle -- anciennnement désignée sous le nom de CujasNum-- sous l'intitulé \enquote{Cours de droit Carbonnier}. Cette entreprise a représenté un chantier d'une rare intensité pour l'éblissement. En effet, les documents ont été numérisé \enquote{en mode texte, avec une relecture intégrale}. Soucieuse d'offrir le \enquote{meilleur accès aux informations contenues dans le texte}, la bibliothèque Cujas a pris le parti d'assurer elle-même la correction complète de l'OCR (\textit{Optical Character Recognition}), ses équipes rectifiant manuellement  chaque coquille de reconnaissance afin d'atteindre un taux de précision de 100\%. 

Cette numérisation intense amène à la création d'une \enquote{plate-forme de recherche [...]par l'ingénieur informatique de la bibliothèque; des sommaires dynamiques ont également été créés}. Afin de s'adapter au mieux aux exigences des chercheurs, l'ergonomie de l'interface de recherche a été profondément améliorée. Elle permet désormais de croiser les \enquote{recherches sur les métadonnées, la table des matières et le plein texte, soit en favorisant la sérendipidité grâce à une navigation par listes d'auteurs ou de titres, ou encore à travers un thésaurus spécialement conçu pour la bibliothèque numérique}. Pour parfaire ce dispositif et optimiser le confort de lecture, l'établissement y adjoindra plus tard une visionneuse perfomante, autorisant un accès \enquote{à une page précise sans télécharger tout le document}.

La numérisation de la collection Jean Carbonnier est la première opération de numérisation interne de la bibliothèque Cujas. Elle a donc permis de tester \enquote{en grandeur réelle, la chaîne de numérisation qu'elle mettait en place pour son propre programme}. Face au succès de cette première numérisation, les équipes la réemploient au moment de son deuxième plus gros projet de numérisation, la liste Pfister-Roumy.